\begin{figure}[H]
\centering
\begin{tikzpicture}[
    neuron/.style={circle, draw=blue!50, fill=blue!8, minimum size=0.75cm, font=\scriptsize},
    out/.style={circle, draw=green!60!black, fill=green!8, minimum size=0.75cm, font=\scriptsize},
    arrow/.style={-{Stealth}, thick},
    back/.style={-{Stealth}, thick, dashed, red!70},
    layer/.style={font=\small\bfseries, gray}
]

% Capa de entrada (3 nodos)
\foreach \i/\y in {1/2, 2/0.7, 3/-0.7}{
    \node[neuron] (i\i) at (0, \y) {$x_{\i}$};
}

% Capa oculta 1 (4 nodos)
\foreach \i/\y in {1/2.7, 2/1.2, 3/-0.2, 4/-1.5}{
    \node[neuron] (h1\i) at (2.5, \y) {};
}

% Capa oculta 2 (4 nodos)
\foreach \i/\y in {1/2.7, 2/1.2, 3/-0.2, 4/-1.5}{
    \node[neuron] (h2\i) at (5, \y) {};
}

% Capa de salida (2 nodos)
\foreach \i/\y in {1/1.0, 2/-0.3}{
    \node[out=45, in=135] (o\i) at (7.5, \y) {$y_{\i}$};
}

% Conexiones entrada → oculta1
\foreach \i in {1,2,3}{
    \foreach \j in {1,2,3,4}{
        \draw[-{Stealth}, gray!50, thin] (i\i) -- (h1\j);
    }
}
% Conexiones oculta1 → oculta2
\foreach \i in {1,2,3,4}{
    \foreach \j in {1,2,3,4}{
        \draw[-{Stealth}, gray!50, thin] (h1\i) -- (h2\j);
    }
}
% Conexiones oculta2 → salida
\foreach \i in {1,2,3,4}{
    \foreach \j in {1,2}{
        \draw[-{Stealth}, gray!50, thin] (h2\i) -- (o\j);
    }
}

% Flecha de forward pass (arriba)
\draw[arrow, blue!70] (-0.6, 3.5) -- node[above, font=\small] {Forward pass $\longrightarrow$} (8.1, 3.5);

% Flecha de backward pass (abajo)
\draw[back] (8.1, -2.3) -- node[below, font=\small] {$\longleftarrow$ Backpropagation $\nabla\mathcal{L}$} (-0.6, -2.3);

% Etiquetas de capas
\node[layer] at (0,   -3.3) {Entrada};
\node[layer] at (2.5, -3.3) {Oculta 1};
\node[layer] at (5,   -3.3) {Oculta 2};
\node[layer] at (7.5, -3.3) {Salida};

\end{tikzpicture}
\caption{Perceptrón multicapa (MLP). El \textit{forward pass} propaga la
información de izquierda a derecha. La retropropagación (\textit{backpropagation})
calcula los gradientes de la pérdida $\mathcal{L}$ en sentido inverso,
actualizando los pesos mediante descenso del gradiente.}
\label{fig:mlp}
\end{figure}
